\section{sprite/exec.cpp File Reference}
\label{exec_8cpp}\index{sprite/exec.cpp@{sprite/exec.cpp}}
Implements functions related to execution. 

{\tt \#include \char`\"{}sprite/exec.hpp\char`\"{}}\par
{\tt \#include \char`\"{}sprite/node.hpp\char`\"{}}\par
{\tt \#include \char`\"{}sprite/operators.hpp\char`\"{}}\par
{\tt \#include $<$deque$>$}\par


Include dependency graph for exec.cpp:\begin{figure}[H]
\begin{center}
\leavevmode
\includegraphics[width=325pt]{exec_8cpp__incl}
\end{center}
\end{figure}
\subsection*{Namespaces}
\begin{CompactItemize}
\item 
namespace {\bf sprite}
\end{CompactItemize}
\subsection*{Classes}
\begin{CompactItemize}
\item 
struct {\bf sprite::Pool\-Item}
\begin{CompactList}\small\item\em An item in the computation pool. \item\end{CompactList}\item 
struct {\bf sprite::Computation\-Pool}
\begin{CompactList}\small\item\em A computation pool. \item\end{CompactList}\end{CompactItemize}
\subsection*{Functions}
\begin{CompactItemize}
\item 
void {\bf execute} (Program const \&pgm, Node \&goal, Yield\-Handler const \&out, Trace\-Option trace)
\begin{CompactList}\small\item\em Execute a goal in the context of the given program. \item\end{CompactList}\item 
void {\bf print\_\-node} (Node const \&node)
\begin{CompactList}\small\item\em Print an expression. \item\end{CompactList}\end{CompactItemize}
\subsection*{Variables}
\begin{CompactItemize}
\item 
Program const $\ast$ {\bf g\_\-program} = NULL\label{namespacesprite_a22}

\begin{CompactList}\small\item\em Global pointer to the currently executing program. \item\end{CompactList}\end{CompactItemize}


\subsection{Detailed Description}
Implements functions related to execution. 



Definition in file {\bf exec.cpp}.

\subsection{Function Documentation}
\index{exec.cpp@{exec.cpp}!execute@{execute}}
\index{execute@{execute}!exec.cpp@{exec.cpp}}
\subsubsection{\setlength{\rightskip}{0pt plus 5cm}void sprite::execute (Program const \& {\em program}, Node \& {\em goal}, Yield\-Handler const \& {\em handler} = {\tt YieldHandler\_\-$<$\ void\ $>$()}, Trace\-Option {\em trace} = {\tt NO\_\-TRACE})}\label{namespacesprite_file_a50}


Execute a goal in the context of the given program. 

This is the the primary entry point to the execution system. Execution proceeds in-place, meaning that the goal expression will likely be modified.

\begin{Desc}
\item[Parameters:]
\begin{description}
\item[{\em program}]The program definition. This provides the necessary information to carry out execution as a series of rewrite steps. \item[{\em goal}]The goal statement. Modified in place. \item[{\em handler}](optional) Defaults to a handler that discards the generated results. The object that handles output as it is generated. \item[{\em trace}](optional) Defaults to true. Enables debug tracing. \end{description}
\end{Desc}


Definition at line 102 of file exec.cpp.

References sprite::Fingerprint::at(), sprite::fair\_\-normalize(), sprite::g\_\-program, sprite::Fingerprint::has(), sprite::Node::id(), sprite::is\_\-norm(), sprite::Fingerprint::resize(), sprite::Fingerprint::set(), sprite::Fingerprint::size(), sprite::Node::tag(), and sprite::Yield\-Handler::yield().

Here is the call graph for this function:\begin{figure}[H]
\begin{center}
\leavevmode
\includegraphics[width=282pt]{namespacesprite_a50_cgraph}
\end{center}
\end{figure}
\index{exec.cpp@{exec.cpp}!print_node@{print\_\-node}}
\index{print_node@{print\_\-node}!exec.cpp@{exec.cpp}}
\subsubsection{\setlength{\rightskip}{0pt plus 5cm}void sprite::print\_\-node (Node const \&)}\label{namespacesprite_file_a51}


Print an expression. 

\begin{Desc}
\item[Note:]Alias for operator$<$$<$, callable from the debugger. \end{Desc}


Definition at line 189 of file exec.cpp.